\documentclass[11pt, a4paper]{article}

\usepackage[T1]{fontenc}
\usepackage[utf8]{inputenc}
\usepackage{tgpagella}
\usepackage[spanish, es-noshorthands]{babel}
\usepackage{amsmath, amssymb}
\usepackage{booktabs}
\usepackage[most]{tcolorbox}
\usepackage[american]{circuitikz}
\usepackage[margin=2.5cm]{geometry}
\usepackage{fancyhdr}

\pagestyle{fancy}
\fancyhf{}
\setlength{\headheight}{16pt}
\lhead{\textbf{UCB Sede Tarija} | Circuitos Electrónicos II}
\chead{Unidad 2 - Op-Amps Ideales}
\rhead{Hoja de Ejercicios}
\renewcommand{\headrulewidth}{0.4pt}
\definecolor{ucbblue}{RGB}{0, 51, 102}

\begin{document}

\begin{center}
    \Large\textbf{\textcolor{ucbblue}{Ejercicios: Amplificadores Operacionales Ideales}}[cite: 6]
\end{center}

\section*{Ejercicio 1 — Reconstrucción: Inversor[cite: 6]}
Sin utilizar inicialmente la fórmula de ganancia predefinida. Considere un amplificador inversor ideal con $R_{in} = 5\,\text{k}\Omega$, $R_f = 15\,\text{k}\Omega$ y $V_{in} = 0.4\,\text{V}$[cite: 6]. Determine:
\begin{enumerate}
    \item La tensión en la entrada no inversora $V_+$[cite: 6].
    \item La tensión en la entrada inversora $V_-$[cite: 6].
    \item La corriente que circula por $R_{in}$[cite: 6].
    \item La tensión de salida $V_o$[cite: 6].
    \item El signo y magnitud de la ganancia $A_v$[cite: 6].
\end{enumerate}

\section*{Ejercicio 2 — No Inversor[cite: 6]}
Se desea diseñar un amplificador no inversor con una ganancia $A_v = 6$ utilizando una resistencia a tierra $R_g = 10\,\text{k}\Omega$[cite: 6]. Determine:
\begin{enumerate}
    \item El valor requerido para la resistencia de realimentación $R_f$[cite: 6].
    \item La tensión de salida $V_o$ para una entrada $V_{in} = 0.25\,\text{V}$[cite: 6].
    \item Por qué la salida de esta configuración no invierte la polaridad de la entrada[cite: 6].
    \item Qué impedancia de entrada predice el modelo ideal para este circuito[cite: 6].
\end{enumerate}

\section*{Ejercicio 3 — Amplificador Sumador[cite: 6]}
Considere un amplificador sumador inversor con $R_f = 30\,\text{k}\Omega$, $R_1 = 10\,\text{k}\Omega$, y $R_2 = 30\,\text{k}\Omega$[cite: 6]. Las tensiones de entrada son $V_1 = 0.1\,\text{V}$ y $V_2 = 0.3\,\text{V}$[cite: 6]. \\
Determine la tensión de salida $V_o$ comenzando el análisis obligatoriamente desde la ecuación de nodos (KCL) en el terminal inversor[cite: 6].

\section*{Ejercicio 4 — Comparación Conceptual[cite: 6]}
Dos estudiantes debaten sobre las características de los Op-Amps: \\
\textbf{Estudiante A:} "El inversor tiene menor impedancia de entrada porque el Op-Amp inversor absorbe corriente."[cite: 6] \\
\textbf{Estudiante B:} "La diferencia de impedancia de entrada proviene exclusivamente de cómo se conecta la fuente a la topología del circuito."[cite: 6] \\
¿Quién tiene razón? Justifique su respuesta basándose en el circuito ideal[cite: 6].

\section*{Ejercicio 5 — Predicción Previa a LTSpice[cite: 6]}
Considere un amplificador no inversor con $R_g = 10\,\text{k}\Omega$, $R_f = 20\,\text{k}\Omega$, excitado por una señal $V_{in}(t) = 0.5\sin(\omega t)\,\text{V}$[cite: 6]. \\
\textbf{Antes de abrir el simulador}, escriba sus predicciones:
\begin{itemize}
    \item $A_v =$ \rule{2cm}{0.4pt}[cite: 6]
    \item $V_{out, p} =$ \rule{2cm}{0.4pt}[cite: 6]
    \item Polaridad = \rule{2cm}{0.4pt}[cite: 6]
\end{itemize}
\textbf{Después de simular}, registre:
\begin{itemize}
    \item $A_{v, \text{sim}} =$ \rule{2cm}{0.4pt}[cite: 6]
\end{itemize}
Explique cualquier discrepancia observada entre la predicción y la simulación[cite: 6].

\end{document}
